\section{Related Work}
\label{sec:related}

\subsection{Trust-based detection and routing}

Trust management in MANETs has developed along two largely separate lines.
The first evaluates node reliability from forwarding behaviour, reputation
and recommendation, and uses the resulting score to steer route
selection~\cite{cai2019evolutionary,xia2013trust,sirajuddin2021tbsmr,singh2020self}.
Approaches in this line combine direct observation with neighbour
recommendation~\cite{singh2020self,sanchez2015model} and, more recently,
with learned models~\cite{dhand2025secure,mahamune2024trust}. The recurring
design questions are how to weight one's own observations against those of
others, how to age evidence, and how to resist a recommender that lies.

Three design axes recur. The first is what the trust score is computed
from. Early work used packet-forwarding ratio and drop
behaviour~\cite{xia2013trust,sanchez2015model}. Later schemes added energy,
mobility and link stability~\cite{veeraiah2021trust,thirunavukkarasu2022cluster},
and recent ones learn the score from features using gradient boosting or
neural models~\cite{ryu2023reputation,dhand2025secure}. The second is what
the score is used for. Most schemes feed it into route selection, either by
penalising low-trust hops in a single-path
protocol~\cite{pathan2018efficient,pari2023enhanced} or by distributing
traffic across trust-weighted multiple
paths~\cite{sirajuddin2021tbsmr,srilakshmi2021improved,srilakshmi2022secure}.
A smaller group treats trust as an admission-control question,
authenticating or excluding nodes before they
participate~\cite{tu2021active,lwin2020blockchain}. The third is how
recommendations from other nodes are weighted and filtered, since a
recommender may itself be
adversarial~\cite{singh2020self,rajeswari2021efficient,pari2023hybrid}.

Our work sits in the second and third of these but differs in what it
optimises. Almost all of the above evaluate the trust mechanism by its
effect on delivery ratio, throughput and delay, which measures the benefit
of excluding attackers. Far fewer measure the cost of excluding honest nodes
by mistake, and fewer still measure how long such an exclusion persists.
Schemes that also address energy and route
stability~\cite{thirunavukkarasu2022cluster,ponguwala2019e2sr} make the cost
side explicit for energy but not for availability. We argue in
Section~\ref{sec:framework} that a cumulative false-isolation measure
belongs alongside the false-positive rate, and Section~\ref{sec:results}
shows the two can differ by more than an order of magnitude for the same
detector.

A consistent finding across this line is that static thresholds produce
false positives under changing topology, which motivates dynamic or adaptive
trust state. Chatzidakis and Hadjiefthymiades~\cite{chatzidakis2022trust}
make this explicit, proposing a mechanism that identifies malicious nodes,
alerts the network, and subsequently restores trust when a node resumes
healthy behaviour or when an earlier classification proves incorrect. The
possibility of restoration is important and we adopt it: a verification
scheme that can only ever remove nodes accumulates errors monotonically, a
point our clean-network experiment makes quantitatively in
Section~\ref{sec:phi}.

\subsection{Learned trust models}

A recent line of work replaces hand-specified trust arithmetic with learned
models. Gradient-boosted trees have been applied to trust-value estimation
for malicious node detection~\cite{dhand2025secure}, feature-optimisation and
classification pipelines to intruder isolation~\cite{kavitha2019india}, and
reinforcement learning to reputation-driven route
selection~\cite{ryu2023reputation}. The attraction is that the decision
boundary is fitted rather than assumed, which matters when the relationship
between evidence and misbehaviour is not monotone or not separable by a
threshold.

The coverage failure reported in Section~\ref{sec:results} is not, however,
a modelling failure. Under the false-location attack the delivered-throughput
trace is identical to the clean-network trace in every seed, so the
behavioural evidence stream contains no information about the attack. No
estimator recovers signal from an identical trace, and a learned model
restricted to that evidence class would fail exactly as the statistical one
does. Evidence coverage is a property of what is collected, not of how it is
combined. A learned rule could contribute at the combination stage; it
cannot address a limitation upstream of it.

\subsection{Localisation-centred beacon security}

The second line is beacon-specific, using ranging measurements and geometric
consistency to determine whether an advertised position is
plausible~\cite{kuriakose2023embn,liu2008attack}. Robustness here usually
means tolerating a bounded number of lying beacons rather than identifying
them individually, and the classical results bound achievable localisation
error as a function of the fraction of compromised anchors.

Secure localisation has a longer history in sensor networks. SeRLoc bounds a
node's position using directional beacons without range
measurements~\cite{lazos2004serloc}; verifiable multilateration rejects
position claims inconsistent with distance bounds~\cite{capkun2006secure};
and Liu, Ning and Du detect malicious beacons by the inconsistency between
measured and declared range, revoking a beacon once a base station has
received enough independent alerts~\cite{liu2005detecting}. Our localisation
branch uses the same residual as the last of these. What differs is the
setting --- mobile beacons and commodity 802.11 RSSI rather than static
beacons or UWB ranging --- and the decision rule: our confirmation is
reversible, and we measure what a false confirmation costs while it
persists.

The two lines overlap less than their subject matter suggests, and the gap
is visible in the evaluations rather than in the methods: a trust-routing
paper does not typically test an attacker whose only misbehaviour is
positional, and a beacon-security paper does not typically test one whose
only misbehaviour is in forwarding. This is an observation about the papers
we surveyed, not a systematic survey finding.

\subsection{EMBN-MANET}

EMBN-MANET~\cite{kuriakose2023embn} is the recent work closest to ours. Its
detection evidence consists of RSSI measurements and trilateration in a
UWB-based mobile ad hoc network. A $k$-polytopes procedure identifies
centroids, from which an identifier infers suspicious activity, and a valid
probable location $T_{loc}$ is compared against the location $F_{loc}$ a
suspect beacon shares during coordinated localisation. Its threat model is a
beacon that masquerades or spoofs by providing false location coordinates.
Its response is to assign a reputation value to the identified beacon and to
exclude it from coordinating localisation once that value crosses a
threshold. Reported identification is complete when fewer than four
malicious beacons occupy a neighbourhood and exceeds 85\% above that count.
The authors state that their work addresses static beacon nodes and name
detection for mobile beacons as future work.

Four differences follow. First, EMBN-MANET's evidence is entirely
localisation-derived. A beacon that advertises correct coordinates while
selectively dropping transit traffic falls outside what its measurements can
observe, and outside the threat model it sets out to address. Second, its
response excludes a beacon from localisation coordination. It does not
remove the node from packet forwarding, so it makes no claim about
post-removal route continuity or throughput, and it is not a model of the
forwarding-level isolation we evaluate. None of our isolation policies is an
implementation of its response. Third, its beacons are static and ours are
mobile, which its authors identify as the natural extension. Fourth, its
accuracy figure is conditional on attacker density within a neighbourhood,
and we quote it as such.

\begin{table}[!t]
\caption{Evidence classes and observable attack surface by approach family.
Entries are read from published descriptions, not measured.}
\label{tab:families}
\centering
\small
\begin{tabular}{@{}lccccc@{}}
\toprule
Approach & Loc.\ evid. & Beh.\ evid. & Rec.\ trust & Mob.\ beac. & FI metric \\
\midrule
Localisation-centred~\cite{kuriakose2023embn,liu2005detecting,liu2008attack} & \checkmark & & & & \\
Behavioural trust~\cite{xia2013trust,sanchez2015model} & & \checkmark & \checkmark & \checkmark & \\
Multipath trust~\cite{sirajuddin2021tbsmr,srilakshmi2021improved} & & \checkmark & \checkmark & \checkmark & \\
Learned trust~\cite{ryu2023reputation,dhand2025secure} & & \checkmark & & \checkmark & \\
Admission control~\cite{tu2021active,lwin2020blockchain} & & \checkmark & \checkmark & \checkmark & \\
This work & \checkmark & \checkmark & \checkmark & \checkmark & \checkmark \\
\bottomrule
\end{tabular}
\par\smallskip\footnotesize
Loc./Beh.\ evid.: localisation and behavioural evidence. Rec.\ trust:
recommendation-based trust. Mob.\ beac.: mobile beacons. FI metric: reports
a false-isolation cost measure.
\end{table}

\subsection{Structured comparison and basis for comparison}

Because we decline a numerical accuracy comparison, we give a structured
qualitative one instead. Table~\ref{tab:families} states, for each family of
approach, which evidence classes it observes and which attack classes
therefore fall inside its scope. Localisation-centred beacon security
observes the evidence class that the false-location attack disturbs and not
the class that selective forwarding disturbs; behavioural trust schemes
observe the converse. Our contribution is not a better detector within
either column but the observation that the columns are disjoint and that a
beacon can misbehave in whichever one is not being watched. The final column
is the metric introduced in Section~\ref{sec:framework}; we are not aware of
prior work in this area that reports the persistence of false exclusions
rather than their rate, though we offer that as a gap in our reading, not an
exhaustive search.

Reported detection accuracy in this field is conditioned on an attack model,
a topology, a mobility regime, a traffic load and an attacker density, all
of which vary across papers. It also saturates: when published figures
cluster between 90\% and 99\%, a difference of one percentage point carries
less information than the differences in evaluation conditions that
produced it. Our own results show why. Every configuration tested, including
both single-evidence ones, achieves perfect recall against the attack its
evidence can observe and near zero against the attack it cannot. An accuracy
figure drawn from the first case alone would be correct and uninformative.
The informative axis is coverage, and we compare against EMBN-MANET on it,
where the comparison is both fair and decidable.
